\documentclass[runningheads,orivec]{llncs}
\usepackage[T1]{fontenc}
\usepackage[utf8]{inputenc}
\usepackage[english]{babel}
\usepackage{amsmath,amssymb}
\usepackage{stmaryrd}
\usepackage{graphicx}
\usepackage{booktabs,tabularx,array}
\usepackage{xcolor}
\usepackage{xspace}
\usepackage{listings}
\usepackage{microtype}
\usepackage[hidelinks]{hyperref}
\urlstyle{same}
\newcommand{\barrel}{\textsc{BARReL}\xspace}
\newcommand{\WD}{\mathsf{WD}}
\newcommand{\enc}[1]{\llbracket #1\rrbracket}
% B's partial-function-space arrow, as in the original manuscript.
\newcommand\pfun{\mathrel{\ooalign{\hfil$\mapstochar\mkern5mu$\hfil\cr$\to$\cr}}}
\lstdefinelanguage{LeanSnippet}{
  morekeywords={import,machine,refinement,implementation,pog,from,obligation,next,of,by,qed,theorem,def,exact,open},
  sensitive=true,
  morecomment=[l]{--},
  morestring=[b]"
}
\lstset{language=LeanSnippet,basicstyle=\small\ttfamily,
  literate={⟨}{{$\langle$}}1 {⟩}{{$\rangle$}}1,
  keywordstyle=\bfseries,commentstyle=\itshape,
  columns=fullflexible,keepspaces=true,showstringspaces=false,
  breaklines=true,frame=single,framerule=0.3pt,
  aboveskip=0.8em,belowskip=0.8em}


% Target: iFS 2027 regular research paper, double-blind submission.
\newif\ifanonymous
\anonymoustrue
\newif\ifdraftnotice
\draftnoticetrue

\title{\barrel: A modern backend for Atelier B in Lean}
\titlerunning{BARReL: A modern backend for Atelier B in Lean}
\ifanonymous
  \author{Anonymous authors}
  \authorrunning{Anonymous authors}
  \institute{Anonymous institution}
\else
  \author{Ghilain Bergeron\orcidID{0009-0008-2814-3480}
    \and Vincent Tr\'elat\orcidID{0009-0006-4143-3939}}
  \authorrunning{G. Bergeron and V. Tr\'elat}
  \institute{Universit\'e de Lorraine, CNRS, Inria, LORIA, Nancy, France}
\fi
\hypersetup{pdftitle={BARReL: A modern backend for Atelier B in Lean},pdfauthor={}}

\begin{document}
\maketitle
\begin{abstract}
\barrel is a Lean~4 backend for Atelier~B that allows users to prove
B proof obligations interactively while retaining their B specifications.
It translates obligations generated by Atelier~B into Lean propositions
over a set-based mathematical library, with notation and lemmas for
reasoning about B sets, relations, and functions. Extrema, function
application, sequence size, and cardinality are represented by operators
whose types require explicit evidence of well-definedness. The translation
generates these conditions in their local logical context and shares
repeated conditions within an import. Extensible automation discharges
routine goals, while users prove the remaining obligations with ordinary
Lean tactics and reusable lemmas. The resulting proof terms are checked
by Lean's kernel, and a completion command publishes the imported
development's declarations. We describe the representation, the handling
of partiality, and the automatic and interactive proof workflow. We
validate the approach on a sensor-monitoring machine whose partial
readings, quorum counts, and guarded extrema support an authorization
invariant. The complete development discharges 80 source obligations
and 66 generated WD goals. Of the latter, 51 are discharged automatically
with reusable case-study lemmas; all remaining proofs are completed
through the interactive interface. The published declarations contain
no admitted proofs or nonstandard axioms.
\keywords{B method \and Lean~4 \and Well-definedness \and Interactive theorem proving}
\end{abstract}

\ifdraftnotice
\begin{center}
\small\itshape Working draft --- 30 September 2026.
\end{center}
\fi

\section{Introduction}
\label{sec:introduction}

The B method~\cite{AbrialBBook} supports the development of software from
mathematical specifications through successive refinements. Its industrial
toolsets, notably Atelier~B~\cite{AtelierB}, generate proof obligations
expressing properties such as invariant preservation and refinement
correctness. Users discharge these obligations with automatic provers
and interactive proof. The quality of this interaction matters: a proof
backend should expose understandable statements, provide reusable
mathematics, and make the assumptions behind each proof step explicit.

Lean~4~\cite{lean} offers a complementary environment for this work. Its
kernel checks proof terms independently of the tactics that construct
them, and its metaprogramming framework supports domain-specific proof
interfaces. Mathlib~\cite{mathlib4} provides an extensive mathematical
library. Bringing these facilities to existing B developments requires
an interpretation of B's set-based expressions in Lean's typed logic and
a treatment of the partial operations that occur throughout B proofs.

Partiality is central to this translation. A function application requires
an argument in the function's domain; a cardinality requires a finite
set; a minimum requires a set with a least element. These conditions
depend on the logical context of the expression. An operation guard may
justify a partial application that would be meaningless elsewhere.
Moreover, an interactive proof can introduce additional terms whose
well-definedness must also be justified. Generating a separate collection
of side conditions is useful only if the subsequent reasoning respects
their connection to the terms they justify.

This paper presents \barrel (B Automated tRanslation for Reasoning
in Lean), a backend that connects Atelier~B proof obligations to Lean's
interactive proof environment. Users retain their B source files.
\barrel invokes Atelier~B when generation is required, reads the exported
obligations, and constructs propositions over a library of B operators.
Sets and relations remain available as mathematical objects rather than
being replaced by an opaque verification condition language. Users can
therefore combine B-specific lemmas with ordinary Lean tactics and Mathlib.

For its proof-carrying operators, \barrel represents well-definedness as
an argument of the operator itself. During translation, a missing argument
becomes an explicit proof obligation closed over the surrounding variables
and hypotheses. Once proved, that evidence is used in the translated
expression. Repeated conditions can share a theorem, and a collection of
extensible tactics attempts to discharge routine conditions automatically.
Each remaining principal or WD obligation can be proved separately, with
earlier facts available as lemmas. A final command checks completion and
publishes the resulting declarations.

We use a fresh sensor-monitoring machine to make the intended benefit
concrete. A zone is authorized only when sufficiently many sensors have
reported values and their readings lie within configured bounds. The
property combines partial maps, cardinalities, and extrema. Routine
well-definedness reasoning establishes that these expressions have a
meaning; the principal proofs establish that updates preserve the
meaningful authorization property. In particular, a batch of changed
readings must invalidate affected zones while leaving unaffected zones'
authorizations justified. This separation gives a direct way to assess
whether the backend handles WD bookkeeping while leaving the developer
to address the substantive argument.

The paper presents the following contributions:
\begin{itemize}
  \item an integrated Atelier~B-to-Lean backend and mathematical library
    for conducting B proofs with ordinary Lean tactics;
  \item a contextual encoding of supported partial operators with explicit
    proof dependencies, shared WD conditions, and extensible automation;
  \item an interactive workflow for proving and reusing individual
    obligations, with an explicit completion and publication boundary.
\end{itemize}
The case study evaluates the resulting framework through a complete
development. Its 80 source obligations generate 754 WD proof requests,
which sharing reduces to 66 retained conditions. With the case study's
reusable lemmas, 51 of these conditions are discharged automatically;
the remaining obligations are proved through the individual proof
commands. A transitive axiom audit checks every published declaration.
Industrial import experience is supplementary to this focused account
of proof construction and completion.

The account distinguishes the properties checked by Lean from the
connection to the source B development. Kernel checking establishes the
translated propositions under their assumptions; it does not by itself
verify the external obligation generator or the translation. We state
this boundary alongside the implemented operator coverage and the proof
completion discipline.

\section{Representing B Obligations in Lean}
\label{sec:encoding}

\subsection{B developments and the translation pipeline}

A B machine declares sets and constants, state variables, an invariant,
an initialization, and operations. Its proof obligations establish, among
other properties, that initialization satisfies the invariant and that
each operation preserves it under its precondition. Refinement introduces
more concrete state together with a gluing invariant and corresponding
proof obligations. \barrel accepts machines, refinements, implementations,
or an already generated POG file. It invokes Atelier~B's BXML and
proof-obligation generators when generation is needed.

The POG XML format~\cite{POGXML} records mathematical obligations, their
hypotheses, definitions, and typing information. \barrel parses this
format, extracts the context of each goal, and constructs a Lean
expression. For variables $x_1,\ldots,x_n$, hypotheses
$H_1,\ldots,H_m$, and conclusion $G$, the target has the schematic form
\[
  \forall x_1\cdots x_n,\quad
  \enc{H_1}\to\cdots\to\enc{H_m}\to\enc{G}.
\]
The translation also constructs the proof arguments required by partial
operators, described below. Unsupported expressions are reported as
encoding failures and prevent completion of the import.

\subsection{Sets, relations, and functions}

B's mathematical language uses sets, maplets, relational operators, and
quantified predicates. A relation $R$ from $A$ to $B$ is a subset of
$A\times B$. A partial function is a functional relation: whenever
$(x,y)\in R$ and $(x,z)\in R$, we have $y=z$. B writes
$R\in A\pfun B$ for membership in this partial-function space. Totality additionally
requires its domain to equal $A$. Thus a B function can also occur as a
relation in domain restriction, image, override, and other set operations.

\barrel uses Mathlib's \texttt{Set} and \texttt{SetRel} representations:
\[
  \mathsf{Set}\ \alpha=\alpha\to\mathsf{Prop},\qquad
  \mathsf{SetRel}\ \alpha\ \beta=\mathsf{Set}(\alpha\times\beta).
\]
B typing information determines the corresponding element types;
integers use Lean's \texttt{Int}, products use product types, and
powersets use \texttt{Set} over the element type. Function properties
are predicates on relations. This lets the same translated object be
used for both function application and relational reasoning.

In the current reader, unknown carrier-type identifiers are mapped to
integers, and deferred carrier sets receive hypotheses stating that
they are finite nonempty integer sets. The case study's sensor and zone
carriers are consequently represented by constrained integer sets,
rather than arbitrary Lean type parameters. The operator library also
provides finite sets, intervals, relational images, extrema, cardinality,
and application, with B-style notation and lemmas for interactive proof.
Operator coverage is an implementation property; acceptance of some
set-based constructs does not imply support for the entire B language.

\subsection{Partial operators with explicit evidence}

A B expression such as $\min(S)$ requires $S$ to have a least element.
Translating it into a total function on all Lean sets would leave a value
available even outside this intended domain. For extrema, application,
sequence size, and cardinality, \barrel instead includes a proof of the
precondition in the operator's type. For minimum on a partially ordered
carrier, the predicate and operator type are
\begin{align*}
  \WD_{\min}(S)&\;\equiv\;
    \exists x\in S,\ \forall y\in S,\ x\leq y,\\
  \mathsf{bmin}&:(S:\mathsf{Set}\ \alpha)\to
    \WD_{\min}(S)\to\alpha.
\end{align*}
The implementation uses the witnessed least element and proves membership
and minimality lemmas. On finite nonempty integer sets, the condition
follows from standard order properties. On an arbitrary ordered carrier,
nonemptiness alone is insufficient. Cardinality requires finiteness;
function application requires the evidence expressed by the library's
application predicate.

An operator occurrence therefore has an explicit proof dependency. During
translation, the encoder creates a fresh metavariable for a missing
proof argument and inserts it into the expression. It then collects the
remaining proof metavariables as WD obligations. Each condition is closed
over the variables and hypotheses available at that occurrence. Once the
condition is proved, the corresponding theorem supplies the expression's
proof argument. Users invoking the same library operator in subsequent
interactive proofs must also supply its evidence.

This discipline is not yet uniform across arithmetic operators. Division,
modulo, and exponentiation currently use direct Lean operators without
the same generated proof arguments. The paper's WD account concerns the
listed proof-carrying operators; the case study uses those
operators and ordinary integer comparisons. The trust boundary of the
translation is discussed in Section~\ref{sec:workflow}.

\subsection{Logical scope of WD evidence}

The position of a partial expression determines which assumptions may
justify it. For example, a reading $r(s)$ can occur after a guard
$s\in\operatorname{dom}(r)$, inside a quantifier whose bound variable is
$s$. Neither the guard nor the variable is a global assumption. The
encoder uses Lean's local context to retain precisely this scope while
constructing the expression and its missing proof arguments.

Write $\Gamma$ for the current context and $\enc{P}_\Gamma$ for a
predicate translated in that context. Implication and conjunction are
constructed as follows:
\begin{align*}
  \enc{P\Rightarrow Q}_\Gamma
    &=\forall h:\enc{P}_\Gamma,\ \enc{Q}_{\Gamma,h},\\
  \enc{P\land Q}_\Gamma
    &=\mathsf{DepAnd}\bigl(\lambda h:\enc{P}_\Gamma.\,
      \enc{Q}_{\Gamma,h}\bigr).
\end{align*}
Here $\mathsf{DepAnd}(Q)$ is defined as $\exists h:P,\ Q(h)$, for
$Q:P\to\mathsf{Prop}$. If $Q$ does not depend on $h$, this is equivalent
to ordinary conjunction. The dependent form allows the right conjunct
to contain operators whose evidence uses the left conjunct. The encoder
first translates $P$, then introduces $h$, and only then translates $Q$.
Consequently, a WD condition created within $P$ cannot use $P$ itself;
one created within $Q$ can. This order matters to the generated
conditions, even when reordering the predicates would preserve their
truth whenever all expressions are defined.
Before encoding, normalization splits conjunctive antecedents into
successive implications, preserving their order and exposing their
components as individual assumptions.

For both universal and existential quantification, the encoder introduces
the bound variable with its translated element type before visiting the
body. B set-membership restrictions remain predicates: introducing an
integer variable alone does not establish membership in a sensor set.
A condition created in the body is closed universally over its local
variables and proof assumptions, regardless of whether the source
quantifier is universal or existential. Thus an existential expression
does not by itself supply a particular witness for proving definedness.
For example, under the outer assumption $r\in S\pfun D$, the application in
\[
  \exists s,\quad s\in\operatorname{dom}(r)\land r(s)\le c
\]
requires a condition of the form
\[
  \forall s,\quad s\in\operatorname{dom}(r)
  \to\WD_{\mathrm{app}}(r,s),
\]
with the outer variables and assumptions also quantified. It does not
require that a suitable reading below $c$ exists: that remains the
principal existential goal. Disjunction and biconditional are encoded
with ordinary Lean connectives and introduce no corresponding local
hypothesis for either operand.

\subsection{A contextual WD obligation}

Consider a partial reading map $r$ and a set $A$ of available sensors
belonging to one zone. A lower-bound check uses $\min(r[A])$, where
$r[A]$ is the relational image of $A$. Its WD obligation must establish
that this image has a least element. If the surrounding context contains
\[
  r\in\mathit{SENSOR}\pfun\mathbb Z,
  \qquad A\in\operatorname{FIN}_1(\operatorname{dom}(r)),
\]
then the image is finite and nonempty. Here $\operatorname{FIN}_1(X)$ is
the collection of finite nonempty subsets of $X$. A library lemma proves
\[
  \bigl(r\in\mathit{SENSOR}\pfun\mathbb Z\bigr)
  \to\bigl(A\in\operatorname{FIN}_1(\operatorname{dom}(r))\bigr)
  \to\WD_{\min}(r[A]).
\]
The encoder's actual closed goal additionally quantifies the variables
and retains the hypotheses of the source occurrence. The proof may
establish the finite-nonempty premise from these hypotheses before
applying the lemma.

The context is essential. A summary operation guarded by $A\ne\varnothing$
can justify the minimum; an unguarded summary on empty readings cannot.
The encoding of implication and dependent conjunction allows subsequent
expressions to use proofs of preceding assumptions. The generated WD
statement records those dependencies, rather than merely asking whether
the syntactic expression $\min(r[A])$ is meaningful in isolation.

The resulting WD theorem also does not prove the principal property.
For example, proving that $\min(r[A])$ exists does not show that it exceeds
a configured threshold. The former belongs to definedness reasoning;
the latter is part of the authorization condition and its preservation
proof. This distinction structures both the automation and the case study.

\paragraph{From local evidence to named obligations.}
Let a missing argument have type $W$ in a local context
$x_1:T_1,\ldots,x_n:T_n$, where some $T_i$ are propositions. Its named
obligation has the closed type
\[
  d_W:\forall x_1:T_1,\ldots,\forall x_n:T_n,\ W.
\]
The generated expression retains a reference to this declaration as its
evidence, instantiated in the context of the occurrence. Nested partial
operators may also make one WD declaration depend on another. The importer
therefore instantiates all remaining proof metavariables before retaining
the generated statements. A principal goal mentioning an unproved WD
declaration remains pending: the individual proof command reports that
dependency and requires its proof first. The principal statement and the
WD statements thus form one proof development, rather than unrelated
lists of assertions.

This describes the implementation's evidence discipline for the supported
operators. Once every referenced declaration has a proof, Lean checks the
resulting terms against their types. It does not follow from that check
alone that the translation preserves B semantics; no mechanically verified
correctness theorem for the reader and encoder is claimed here.

\subsection{Sharing repeated conditions}

Different obligations can contain the same closed WD proposition.
\barrel maintains a table of conditions already encountered in the
current import, including both proved and pending ones. Canonicalization
removes binder names and expression metadata; a hash selects a bucket,
and comparison uses the canonical form with a bucket-local definitional
equality fallback. A match reuses the first condition's declaration
instead of creating another proof task.

The context remains part of the key. Conditions justified by different
hypotheses need not be shared even if they mention the same expression.
The procedure recognizes repeated representations, rather than proving
logical equivalence or subsumption. It does not share across imports,
and its hash-bucket restriction is not a complete test of definitional
equivalence. When a shared condition is still pending, its dependent
obligations continue to require its proof. Sharing removes duplicated
work without removing the common precondition.

\section{Automatic and Interactive Proof}
\label{sec:workflow}

\subsection{A library for B reasoning}

The encoding is accompanied by definitions and proved lemmas for the
represented B operators. For example, the domain of a total function is
its source set; a function's image of a finite set is finite; and the
minimum of a finite nonempty image belongs to that image. Such results
serve both interactive proofs and automatic WD discharge. Since sets
and relations use Mathlib representations, users can also apply ordinary
set extensionality, image laws, arithmetic reasoning, and rewriting.

WD automation is extensible through labelled collections of lemmas.
The implementation provides \texttt{wd\_min}, \texttt{wd\_max},
\texttt{wd\_app}, and \texttt{wd\_card} collections and corresponding
tactics. A theorem registered with a label becomes available to the
associated search. For a linearly ordered codomain, such as the integers
used in the case study, the minimum collection contains a rule of the form
\[
  f\in A\pfun B\quad\land\quad
  S\in\operatorname{FIN}_1(\operatorname{dom}(f))
  \quad\Longrightarrow\quad \WD_{\min}(f[S]),
\]
where $\operatorname{FIN}_1(X)$ denotes the finite nonempty subsets of $X$.
The same pattern supports maximum; cardinality uses finiteness rules;
application uses function properties and domain membership.

The user-facing \texttt{barrel\_solve} tactic combines basic introduction
and triviality steps, B typing tactics, specialized WD reasoning, and
Lean's \texttt{grind} tactic. The specialized tactics introduce assumptions,
substitute equalities, and generalize proof arguments before searching
for an applicable lemma. This reduces sensitivity to the particular
proof term supplied as a WD argument. All successful attempts construct
Lean proofs; a failed attempt leaves a residual obligation for the user.

A domain-specific lemma can connect an invariant to a reusable WD rule.
In the monitoring case study, a helper proves that a permitted zone's
available sensors form a finite nonempty subset of the reading map's domain.
The general finite-image rules then establish the extrema conditions.
Section~\ref{sec:evaluation} reports the automatic discharge rate with
these helpers and the remaining explicit WD proofs. The measured
configuration includes this extension to the distributed automation.

\subsection{Individual obligations and reusable proofs}

Each import creates a collection of principal and WD obligations and
attempts automatic discharge. Users prove the remaining goals by name
or request the next pending obligation. The following excerpt is taken
from the completed case study. The ellipses omit other proved obligations;
the complete source contains no proof placeholders:
\par\noindent\begin{minipage}{\linewidth}
\begin{lstlisting}
import pog ZoneMonitor from "specs/zonemonitor"
-- ... other obligation commands ...
obligation Initialisation_0 of ZoneMonitor by
  intros
  exact ⟨Set.empty_subset _, by simp⟩
-- ... other obligation commands ...
qed ZoneMonitor
\end{lstlisting}
\end{minipage}\par
The \texttt{by} form accepts a tactic script, and \texttt{from} accepts a
proof term. The \texttt{of} clause selects an import; without it, the
latest import is used. Names are relative to the selected import,
including relative names of WD obligations.

A completed obligation is available to subsequent proofs in the same
import before final publication. A user can therefore establish a
function-image lemma or a common invariant consequence once and use it
in several proofs. Ordinary Lean definitions and lemmas can also be
introduced between obligation commands. Imports may be interleaved,
with each retaining its own progress and proof state. Diagnostics and
proof skeletons help locate residual goals and distinguish automatic
proofs, interactive proofs, admissions, and unsuccessful encodings.

Separate commands also allow Lean to retain unchanged command prefixes
during editing. Editing a late proof need not replay earlier proofs;
editing an earlier command causes the affected suffix to be elaborated
again. Tactics execute synchronously. These are properties of the proof
interface and its implementation, without implying automatic transport
of proofs after a change to the B specification.

\subsection{Completion and implementation}

An import retains a base and a working Lean environment. The latter
contains its generated declarations and completed proofs, which are
available while proving that import but remain unpublished in the
surrounding environment. The \texttt{qed} command checks that no obligation
is pending and no source goal failed to encode, then merges the checked
declarations into the current global environment. The merge rejects
name collisions, verifies replayed declarations against the kernel
environment, and restores metadata needed by generated definitions.
Nothing is published until that merge succeeds.

To avoid unnecessary replay, the implementation reuses a checked working
environment when a conservative comparison detects no external change.
New declarations or changed attributes can trigger a checked rebase.
Indexed obligation names, a pending-goal cursor, and maintained progress
counts support selection and reporting. These details implement the
same logical workflow for imports containing many obligations.

Automatic proof attempts receive individual heartbeat budgets, and
unsuccessful attempts leave their goals pending. Encoding failures are
recorded and block completion. Process-level failures such as exhaustion
of memory remain possible; successful import is consequently a separate
observation from complete discharge of the imported development.

\subsection{Trust boundary and proof completion}

Lean checks each proof against its translated statement and assumptions.
The connection to B additionally relies on Atelier~B's generation of the
obligations, the XML reader, extraction, and the encoding. A valid proof
of an incorrectly translated proposition would not establish the intended
source property. \barrel therefore provides kernel-checked proofs of
its generated statements, rather than a verified end-to-end translation.

The completion command follows Lean's convention for explicit admissions:
a proof containing \texttt{sorry} can be accepted with a warning. For a
completed case-study result, successful \texttt{qed} must therefore be
accompanied by an audit of the published declarations and their
dependencies, excluding admissions and unintended axioms. This separates
proof development in progress from the evidence reported in the paper.

\section{Case Study: A Zone Monitor}
\label{sec:evaluation}

\texttt{ZoneMonitor} is a B machine implemented and imported into \barrel
for this case study. It manages readings from redundant sensors and
maintains authorization for zones whose readings satisfy a configured
policy. Partial-map application, cardinality, and extrema arise from the
policy and its operations. The development exercises the translation of
these expressions, automatic WD discharge, reusable relational lemmas,
and individual proofs of invariant preservation. All 146 retained
obligations are proved and published in Lean without admissions.

\subsection{State and safety policy}

Let $S$ and $Z$ denote the finite nonempty carrier sets \texttt{SENSORS}
and \texttt{ZONES}. A fixed total map $\mathit{zone}\in S\to Z$ assigns
each sensor to a zone, and $q\in Z\to\mathbb{N}_{>0}$ specifies the
required number of available sensors. Measurements lie in a fixed nonempty
integer interval $D$. The state contains
\[
  r\in S\pfun D,\qquad
  \ell,u\in Z\to D,\qquad P\subseteq Z.
\]
Here $r$ contains the available readings, $\ell$ and $u$ give the configured
bounds, and $P$ contains the zones currently authorized. The B variables
are named \texttt{reading}, \texttt{lower}, \texttt{upper}, and
\texttt{permit}, respectively. Define
\[
  A_z=\operatorname{dom}(r)\cap\mathit{zone}^{-1}[\{z\}],
  \qquad V_z=r[A_z].
\]
The quorum counts sensors in $A_z$, not distinct values in $V_z$: two sensors
reporting the same value still constitute two available sensors. The invariant
requires $\ell(z)\le u(z)$ for every $z\in Z$, together with
\begin{equation}
\label{eq:zone-policy}
\forall z\in P,\quad
\begin{aligned}
  &A_z\ne\varnothing\ \land\ q(z)\le\operatorname{card}(A_z)\ \land{}\\
  &\ell(z)\le\min(V_z)\ \land\ \max(V_z)\le u(z).
\end{aligned}
\end{equation}
Non-emptiness precedes the extrema in the invariant. Finiteness of $S$
makes $A_z$ finite, and $A_z\subseteq\operatorname{dom}(r)$ makes its image
nonempty whenever $A_z$ is nonempty. No assumption requires every zone to
have enough installed sensors: a zone with too few can never be authorized.

Initialization sets $r=P=\varnothing$ and assigns initial bound maps
$\ell_0,u_0\in Z\to D$, supplied as constants satisfying
$\ell_0(z)\le u_0(z)$. The policy concerns recorded measurements; it does
not assume that sensors report the true physical values. The resulting
property is authorization consistency within the monitor, not a
certification of the monitored plant.

\subsection{Operations and the batch-update proof}

Table~\ref{tab:zone-operations} summarizes the eight operations. Reading
updates and configuration changes revoke affected authorizations. This
allows readings to change freely while preventing a permit from retaining
a justification based on stale values.

\begin{table}[tb]
\caption{Operations of the implemented \texttt{ZoneMonitor} machine.}
\label{tab:zone-operations}
\small
\begin{tabularx}{\textwidth}{@{}lX@{}}
\toprule
Operation & Guard or state change \\
\midrule
\texttt{Record} & Override $r$ at sensor $s$ and remove $\mathit{zone}(s)$ from $P$. \\
\texttt{RecordBatch} & Override $r$ with $b\in S\pfun D$ and remove $\mathit{zone}[\operatorname{dom}(b)]$ from $P$. \\
\texttt{Withdraw} & Require $s\in\operatorname{dom}(r)$; remove its reading and revoke its zone. \\
\texttt{Configure} & Set the bounds for $z$, requiring the new lower bound not to exceed the upper; revoke $z$. \\
\texttt{Authorize} & Require the conditions in (\ref{eq:zone-policy}) for $z\in Z$; add $z$ to $P$. \\
\texttt{Revoke} & Remove $z\in Z$ from $P$. \\
\texttt{ReadSensor} & Require $s\in\operatorname{dom}(r)$; return $r(s)$. \\
\texttt{ZoneSummary} & Require $z\in Z$ and $A_z\ne\varnothing$; return its cardinality and the minimum and maximum of $V_z$. \\
\bottomrule
\end{tabularx}
\end{table}

The batch operation is written directly in B:
\begin{lstlisting}[language={},morekeywords={PRE,THEN,END}]
RecordBatch(batch) =
  PRE batch : SENSORS +-> VALUE
  THEN
    reading := reading <+ batch ||
    permit := permit - zone[dom(batch)]
  END
\end{lstlisting}
Here \texttt{+->} is the ASCII form of B's partial-function-space arrow
$\pfun$, \texttt{<+} is relational override, and \texttt{||} combines the
two assignments in parallel. Thus both right-hand sides refer to the
old state. The operation updates any subset of sensor readings in one
step, and preserves permits only for zones outside the image of the
updated sensors.

Write $r'=r\mathbin{\oplus}b$ for the new reading relation and
$T=\mathit{zone}[\operatorname{dom}(b)]$. For $z\in P\setminus T$, no
sensor assigned to $z$ belongs to $\operatorname{dom}(b)$. Relational
override leaves both the presence and the value of every reading in that
zone unchanged. The Lean development proves the corresponding frame
lemmas \texttt{available\_overload} and \texttt{readings\_overload}:
\[
  z\notin T\quad\Longrightarrow\quad
  A'_z=A_z\ \land\ r'[A'_z]=r[A_z].
\]
These lemmas are stated for arbitrary relations over arbitrary types;
their proofs use set extensionality, the domain of override, and the
definition of relational image. They do not depend on sensor bounds or
on the authorization invariant. Non-emptiness, quorum, and the two
extremum bounds can then be transported separately to the updated state.

For example, the following is the actual \barrel command proving the
lower-bound preservation obligation:
\begin{lstlisting}[mathescape=true]
obligation Operation_RecordBatch_15 of ZoneMonitor by
  intros
  expose_names
  obtain $\langle$_hne, _hcard, hmin, _hmax$\rangle$ := h_15 zz h_17.1
  simpa only [readings_overload h_17.2] using hmin
\end{lstlisting}
The generated hypothesis \texttt{h\_15} is the old authorization policy;
\texttt{h\_17.1} states that the selected zone was permitted, and
\texttt{h\_17.2} that it is unaffected by the batch. The proof extracts
the old lower bound and rewrites the updated image using the frame lemma.
The maximum and quorum obligations use the same argument. Single-sensor
updates reduce to the singleton case, while withdrawal uses analogous
lemmas for domain subtraction. Configuration changes additionally require
lemmas describing application after a singleton override.

The WD prerequisite of this lower-bound obligation is a separate named
goal. Its proof first rewrites the updated image and then applies the
finite-nonempty-image minimum lemma, using non-emptiness from the old
policy. This order exposes the distinction between establishing that the
minimum exists and proving its lower bound. Both arguments reuse the same
mathematical fact about unaffected zones, but they establish different
propositions.

\subsection{Source WD and generated WD obligations}

The artifact imports a POG generated by Atelier~B with its \texttt{-w}
option. It contains 41 principal obligations and 39 source WD obligations.
The latter include definedness of constants and the invariant, operation
preconditions, and output expressions. This is significant for
\texttt{ReadSensor} and \texttt{ZoneSummary}: these read-only operations
preserve the state, but their application and extrema still require WD
proofs. Their five source WD goals are included in the development.

These source obligations differ from the conditions created by \barrel
when it translates a partial expression. Source WD obligations are B
predicates emitted by Atelier~B. Generated WD obligations are proof
arguments required by the Lean operator definitions. Importing the full
POG therefore proves both families. It also avoids counting the absence
of an invariant-preservation goal for an output-only operation as evidence
that its output expressions are defined. The case uses
\texttt{import pog ZoneMonitor from "specs/zonemonitor"}; the ordinary
machine-import path currently invokes the generator without \texttt{-w}.

Encoding the 80 source goals creates 754 requests for WD evidence.
Sharing retains 66 distinct generated WD goals and reuses 688 earlier
conditions. The repeated requests arise from the same typing hypotheses,
policy clauses, and updated-state expressions across generated goals.
They do not correspond to 754 distinct mathematical facts. The retained
66 conditions, together with the 80 source goals, give 146 obligations
in the \barrel development.

\subsection{Automation and remaining proof work}

The case supplies proved helper lemmas and a bounded WD tactic branch
through \barrel's extension mechanism. The helpers derive finiteness of
available sensors, membership in a map's source, and non-emptiness from
an authorization policy. They then invoke the library's application,
cardinality, and finite-image extrema lemmas. The branch recognizes the
four corresponding WD predicates and uses bounded searches for their
premises. This is the measured configuration; the results are not a claim
about unextended stock automation.

On import, automation discharges 51 of the 66 distinct generated WD goals
($77.3\%$) and 20 of the 80 source goals. The latter comprise 10 principal
and 10 source WD goals. The remaining 75 obligations receive individual
proof commands: 31 principal, 29 source WD, and 15 generated WD goals.
The 29 source WD commands explicitly invoke a reusable
\texttt{zone\_source\_wd} tactic. They are counted as supplied proof
commands, not as successes of import-time automatic discharge. The
remaining generated WD proofs use the frame lemmas and authorization
case distinctions described above.

\begin{table}[tb]
\caption{Obligations by source. Parentheses give import-time automatic
proofs. Generated WD goals are assigned to the source that first introduces
them; later uses contribute requests but no new distinct goal.}
\label{tab:zone-evidence}
\small
\begin{tabular}{@{}lrrrr@{}}
\toprule
Source & Principal & Source WD & WD requests & Unique WD \\
\midrule
Properties & 0 & 4 (4) & 0 & 0 \\
Initialization & 5 (4) & 0 & 16 & 8 (8) \\
Invariant & 0 & 15 (3) & 76 & 4 (4) \\
\texttt{Record} & 6 (1) & 2 (0) & 91 & 15 (13) \\
\texttt{RecordBatch} & 6 (1) & 0 & 66 & 6 (4) \\
\texttt{Withdraw} & 6 (1) & 2 (0) & 91 & 7 (5) \\
\texttt{Configure} & 8 (1) & 0 & 88 & 8 (6) \\
\texttt{Authorize} & 5 (1) & 11 (3) & 220 & 12 (7) \\
\texttt{Revoke} & 5 (1) & 0 & 56 & 6 (4) \\
\texttt{ReadSensor} & 0 & 1 (0) & 10 & 0 \\
\texttt{ZoneSummary} & 0 & 4 (0) & 40 & 0 \\
\midrule
Total & 41 (10) & 39 (10) & 754 & 66 (51) \\
\bottomrule
\end{tabular}
\end{table}

Table~\ref{tab:zone-evidence} reports the counts by source. In particular,
the read-only operations contribute 50 generated requests, all reusing
conditions introduced earlier. Their zero in the final column therefore
indicates sharing, not an absence of definedness reasoning. The evidence
export retains the generated declaration dependencies so that the shared
conditions and their consumers can be inspected.

The completed development uses Lean~4.31.0, with a budget of 20,000
heartbeats per import-time proof attempt and 200,000 for subsequent
proof commands. Running \texttt{qed ZoneMonitor} publishes all 146
declarations. A subsequent audit checks that there are no skipped,
pending, or admitted obligations and examines the transitive axioms of
every published obligation. The only axioms found are Lean's standard
\texttt{propext}, \texttt{Classical.choice}, and \texttt{Quot.sound}.
The artifact retains the B machine, full POG, helper lemmas, individual
proof commands, JSON evidence, tabulated counts, run log, and source
hashes. Thus the counts describe an executed proof development.

The automatic fraction therefore measures routine proof work removed at
import time under a documented extension. The residual proofs identify
where the application policy matters: preserving readings outside an
update, recovering non-emptiness from a permit, and distinguishing an
existing permit from a newly authorized zone. This single machine tests
those interactions without a refinement chain or a separate industrial
benchmark. Industrial inputs can supply additional evidence about import
scale; their size is independent of the proof-completion result for this
case. Kernel checking establishes the translated propositions under the
embedding, with the translation trust boundary stated in
Section~\ref{sec:workflow}.

\section{Related Work}
\label{sec:related}

Bodeveix and Filali study type synthesis and the translation of B to
PVS~\cite{Bodeveix2002BtoPVS}. Their work addresses the passage from B's
set-based language to the typing discipline of the target prover.
\barrel shares this concern and targets obligations exported by Atelier~B,
with an interactive interface hosted in Lean. Its mathematical objects
use Mathlib sets and relations, allowing the imported propositions to be
proved with the host prover's tactics and libraries.

Proof requirements attached to partially specified operations have
established precedents. Sozeau's Russell language adapts predicate
subtyping from PVS to Coq, generating obligations whose proofs complete
a well-typed term~\cite{Sozeau2007SubsetCoercions}. Schmalz studies
partial functions, definitional extensions, and reasoning about the
logic of Event-B, with an Isabelle/HOL integration for
Rodin~\cite{Schmalz2012RodinIsabelle}. These approaches provide relevant
foundations for expressing partiality in a total logic. \barrel applies
explicit proof requirements at occurrences of supported B operators,
closes the resulting conditions over their local context, and connects
them to automatic and interactive discharge of Atelier~B obligations.

The Z Mathematical Toolkit in Isabelle/HOL develops a Z metalanguage
and type universe over HOL, connected to concrete representations such
as finite functions and lists. Coercive subtyping and overloading support
Z-style sets, relations, functions, and sequences~\cite{Foster2025ZToolkit}.
\barrel's library serves the corresponding role for the B fragment
accepted by its translator in Lean. The availability of reusable
mathematical lemmas is important in both settings: the usefulness of
an embedding depends on how users can reason about its representations,
as well as on the expressions it can translate.

Ballenghien and Wolff's HOL-B prototype hosts Event-B as a domain-specific
language in Isabelle/HOL, with a mathematical toolkit, IDE and
documentation support, and generated denotational semantics for machine
specifications~\cite{Wolff2024EB_Isabelle}. This places machine semantics
inside the proof assistant. \barrel retains Atelier~B's proof-obligation
generation and consumes the exported formulas. The resulting workflow
fits existing B developments, while leaving the correctness of the
external generator and of the translation outside the kernel guarantee.
This architectural choice is distinct from interpreting entire machines
inside the prover.

Automated backends complement the interactive approach. D\'eharbe studies
the integration of SMT solvers with B and Event-B environments~\cite{DeharbeBSMT},
and BEer translates B proof obligations into SMT-LIB~\cite{beer}.
\barrel instead exposes propositions and proof tasks in Lean: users may
combine automatic tactics with structured interactive arguments, prove
auxiliary lemmas, and inspect the resulting dependencies. Its automation
is intended to handle routine conditions within that proof workflow.
The case study examines this division of work through the
WD obligations discharged automatically and the principal proofs that
remain, rather than a competition between solver success rates.

\section{Conclusion}
\label{sec:conclusion}

\barrel provides an interactive Lean backend for proof obligations
produced by Atelier~B. Its mathematical library interprets B expressions
using sets and relations, and its supported proof-carrying operators
make the dependencies on well-definedness evidence explicit. The
translation generates contextual WD obligations, shares repeated
conditions, and attempts their discharge with extensible automation.
Users complete the remaining proofs with ordinary Lean tactics and
reusable lemmas before publishing the development's declarations.

The ZoneMonitor machine brings these elements together in one
model: partial sensor readings, quorum counts, and guarded extrema
support an authorization invariant whose preservation also requires
substantive relational reasoning. Its complete BARReL proof discharges
41 principal obligations, 39 source WD obligations, and
66 generated WD goals. Sharing removes 688 repeated WD requests, and
the case study's automation discharges 51 of the 66 retained WD goals.
The remaining proofs use reusable relational lemmas and individual
proof commands. All 146 published obligations pass the transitive
axiom audit without admissions. This case demonstrates the division
between automatic definedness reasoning and explicit policy-preservation
arguments on one machine; it does not establish an automation rate for
arbitrary B developments.

The present guarantees concern the translated Lean propositions.
Verification of the importer, broader operator coverage, and further
integration with B development environments remain directions for
future work.


\subsubsection*{Data Availability Statement}
The accompanying development contains the \barrel sources, pinned
dependencies, the ZoneMonitor B machine, the complete generated POG,
Lean proofs, and scripts for regenerating the obligations and checking
the proof. Machine-readable results record each obligation's status,
automation classification, and transitive axioms, together with source
hashes and the per-operation counts. An anonymized archival review
artifact remains to be deposited.

\bibliographystyle{splncs04}
\bibliography{references}
\end{document}
